\documentclass[11pt]{article}

\usepackage{amsmath}
\usepackage{amssymb}
\usepackage{graphicx}
\usepackage{hyperref}
\usepackage{booktabs}
\usepackage{enumitem}

\title{
EvoMind-DGM: Knowledge-Guided Neuro-Symbolic Evolution\\
for General Intelligence
}

\author{
Yohanes Mesfin \\
iCog Labs / SingularityNET  \\
\text{yohanesmesfin3@gmail.com}
}

\date{}

\begin{document}

\maketitle

\begin{abstract}
Darwin G\"odel Machines (DGMs) demonstrate that language-model-based agents can improve their own software scaffolding through archive-based self-modification and empirical selection. However, existing DGM-style systems are primarily oriented toward task-specific improvement, especially in software engineering. General intelligence requires more than improvement on a narrow benchmark: it requires cumulative knowledge acquisition, abstraction, transfer across domains, reasoning under uncertainty, and attention-guided allocation of limited cognitive resources.

We propose EvoMind-DGM, a knowledge-guided neuro-symbolic architecture for open-ended evolution of self-improving agents. EvoMind-DGM extends the DGM framework with a Hyperon-inspired, MeTTa-centered implementation strategy. A DGM evolution controller orchestrates self-improvement around a symbolic runtime built from current Hyperon-family components, especially MeTTa and PeTTa, together with persistent space-backed memory. This memory layer serves as a knowledge substrate for storing agent lineages, task experiences, learned concepts, symbolic policies, failure modes, transfer patterns, and meta-cognitive hypotheses. MOSES is used to evolve reusable symbolic cognitive policies and procedures, including reflection, verification, memory retrieval, task decomposition, and archive-selection strategies. ECAN provides attention-based resource allocation across agents, tasks, memories, policies, and uncertain hypotheses. PLN, including current PeTTa-based implementations, supplies probabilistic reasoning over causality, transfer, risk, and expected improvement. LLMs remain central as semantic generators and implementation-level mutation operators, but are embedded within a broader neuro-symbolic evolutionary system.

The central hypothesis is that self-improving agents guided by persistent symbolic knowledge, probabilistic reasoning, attention allocation, and symbolic program evolution will produce more interpretable, transferable, and generalizable improvements than DGM systems relying primarily on LLM-generated code modifications and fixed benchmark selection. We outline an experimental framework for evaluating EvoMind-DGM on cross-domain adaptation, knowledge reuse, policy transfer, reasoning-guided mutation, and improvement per compute unit.
\end{abstract}

\section{Introduction}

Recent progress in large language models (LLMs) has enabled increasingly capable software agents that can reason, write code, call tools, inspect files, and modify their own execution environment. Darwin G\"odel Machines (DGMs) extend this direction by allowing agents to iteratively improve themselves through an evolutionary process: a parent agent is selected from an archive, a modified child agent is generated, the child is empirically evaluated, and successful variants are stored for future evolution.

This approach is powerful because it replaces the requirement for formal proof of self-improvement with empirical validation. Rather than proving that a self-modification is beneficial, the system tests whether the resulting agent performs better. This makes DGM-like systems more practical than classical G\"odel Machine formulations, particularly when the agent is implemented as a complex software system around a frozen foundation model.

However, the current DGM framing remains limited as a model of general intelligence. Improving benchmark performance in a single domain, such as software engineering, does not by itself imply broad intelligence. A generally intelligent system should not merely become better at one task class. It should accumulate reusable knowledge over time, form abstractions, reason about its own failures, transfer strategies across domains, allocate attention under limited resources, and revise its beliefs based on experience.

We argue that the path from self-improving agents to general intelligence requires shifting the unit of evolution from task-specific agent code to knowledge-bearing cognitive systems capable of accumulating, abstracting, and transferring experience over time.

To address this, we propose \textbf{EvoMind-DGM}, a neuro-symbolic extension of Darwin G\"odel Machines grounded in a Hyperon-inspired, MeTTa-centered cognitive architecture. EvoMind-DGM combines five core ideas:

\begin{enumerate}[leftmargin=*]
    \item \textbf{DGM-style open-ended evolution}: archive-based self-improvement of agent variants.
    \item \textbf{Space-backed knowledge accumulation}: Atomspace-inspired persistent storage of experiences, concepts, policies, lineages, and meta-cognitive hypotheses.
    \item \textbf{MOSES-based symbolic evolution}: evolutionary discovery of compact, reusable cognitive policies and procedures.
    \item \textbf{ECAN-based attention allocation}: prioritization of agents, memories, tasks, hypotheses, and evolutionary branches.
    \item \textbf{PLN-based probabilistic reasoning}: uncertain inference over causality, transfer, risk, and expected improvement.
\end{enumerate}

In this architecture, LLMs remain essential, but they are not treated as the entire intelligence system. Instead, LLMs act as semantic generators, language interfaces, code/policy proposers, and implementation-level mutation engines within a broader neuro-symbolic evolutionary loop coordinated by a DGM controller and symbolic runtime.

The central claim of this paper is that general intelligence requires not only self-modification, but also persistent knowledge, symbolic abstraction, attention, reasoning, and transfer. EvoMind-DGM is proposed as an architecture for studying these mechanisms in a unified, experimentally testable system.

\section{Background}

\subsection{Darwin G\"odel Machines}

Darwin G\"odel Machines are self-improving agent systems inspired by the original G\"odel Machine concept but grounded in empirical evolutionary selection. A DGM maintains an archive of agent variants. At each iteration, the system selects a parent agent, generates a modified child agent, evaluates the child on a set of tasks, and stores useful variants in the archive.

The basic DGM loop can be represented as:

\begin{equation}
A_{t+1} = \text{ArchiveUpdate}(A_t, \text{Evaluate}(\text{Mutate}(\text{Select}(A_t))))
\end{equation}

where $A_t$ is the archive at time $t$.

In existing DGM systems, the mutation operator is often implemented using an LLM that modifies the agent's code, prompts, tools, or workflow. The evaluation function is typically benchmark-based. This works especially well in software engineering because the ability to modify an agent's software scaffold is itself closely related to the coding ability being evaluated.

However, this coupling between task performance and self-modification ability is domain-specific. General intelligence requires a broader framework in which improvements are not judged only by performance on one class of tasks, but by their ability to support transfer, abstraction, learning efficiency, and robust adaptation.

\subsection{Hyperon, MeTTa, and Space-Based Memory}

Hyperon is a neuro-symbolic AGI framework designed to integrate symbolic reasoning, probabilistic inference, evolutionary learning, attention allocation, and neural/LLM-based methods. In current practice, the Hyperon ecosystem is best understood not as a single monolithic runtime containing every cognitive subsystem, but as a family of interoperable components centered on MeTTa runtimes and symbolic spaces. The practical base includes MeTTa implementations such as \texttt{hyperon-experimental} and PeTTa, together with space APIs and optional backends such as MORK-backed spaces.

In EvoMind-DGM, the long-term cognitive memory of the system is best viewed as an Atomspace-inspired, space-backed symbolic memory layer rather than as a single fixed implementation. Unlike a conventional DGM archive, which primarily stores agent variants and benchmark scores, this memory layer stores both evolutionary and cognitive knowledge:

\begin{itemize}[leftmargin=*]
    \item agent lineages and mutation histories,
    \item task experiences and outcomes,
    \item learned concepts and abstractions,
    \item symbolic cognitive policies,
    \item transfer patterns across domains,
    \item failure modes and safety violations,
    \item probabilistic hypotheses about improvement,
    \item attention values and resource-allocation priorities.
\end{itemize}

This turns the archive from a passive record of variants into an active knowledge base over which the system can reason, regardless of whether the underlying implementation uses Hyperon spaces, PeTTa-managed spaces, or other compatible symbolic-memory backends.

\subsection{MOSES for Symbolic Evolution}

MOSES, or Meta-Optimizing Semantic Evolutionary Search, is an evolutionary program-learning approach that searches over compact symbolic programs or policies. In EvoMind-DGM, MOSES is used not to evolve entire agents at first, but to evolve reusable symbolic cognitive policies.

Examples of policies that MOSES may evolve include:

\begin{itemize}[leftmargin=*]
    \item when to reflect before answering,
    \item when to retrieve memory,
    \item when to call a verifier,
    \item when to use a tool,
    \item when to replan,
    \item when to stop,
    \item how to aggregate evaluation signals,
    \item how to decide whether a child agent enters the archive.
\end{itemize}

This allows EvoMind-DGM to evolve interpretable cognitive control structures rather than relying entirely on opaque LLM-generated modifications. In the current ecosystem, MOSES should be treated as a symbolic subsystem executed through MeTTa-centered workflows, especially repositories such as \texttt{metta-moses}, rather than assumed to be a built-in layer of a single Hyperon runtime.

\subsection{ECAN for Attention Allocation}

ECAN, or Economic Attention Networks, provides a mechanism for allocating limited cognitive resources. In EvoMind-DGM, ECAN is used to assign attention values to agents, tasks, memories, symbolic policies, hypotheses, and evolutionary lineages. In implementation terms, ECAN is best viewed as a symbolic attention subsystem whose outputs are consumed by the DGM controller.

The need for attention allocation arises because an open-ended DGM-style system can produce a rapidly growing archive of variants and experiences. Evaluating every child agent on every task is computationally infeasible. ECAN provides a principled mechanism for determining where compute and reasoning should be focused.

For example, attention may be allocated to:

\begin{itemize}[leftmargin=*]
    \item agents with high transfer potential,
    \item policies that improved multiple domains,
    \item uncertain hypotheses requiring more evidence,
    \item underexplored evolutionary lineages,
    \item tasks likely to reveal generalization failures,
    \item memories relevant to the current mutation objective,
    \item safety risks requiring deeper evaluation.
\end{itemize}

\subsection{PLN for Probabilistic Reasoning}

PLN, or Probabilistic Logic Networks, provides uncertain reasoning over symbolic knowledge. In EvoMind-DGM, PLN reasons over symbolic-memory representations of agents, tasks, policies, failures, and improvements. In the current stack, PLN is naturally treated as a PeTTa/MeTTa reasoning subsystem invoked by the controller rather than assumed to be a built-in service of the base Hyperon runtime.

PLN is used to infer:

\begin{itemize}[leftmargin=*]
    \item whether a mutation likely caused an observed improvement,
    \item whether a policy is likely to transfer across domains,
    \item whether a performance gain is likely benchmark-specific,
    \item whether a child agent introduces safety risk,
    \item which task would best reduce uncertainty,
    \item which evolutionary branch has the highest expected value.
\end{itemize}

In this way, EvoMind-DGM does not merely record which agents performed better. It develops probabilistic beliefs about why they improved and whether those improvements are generally useful.

\section{Motivation: From Task-Specific Improvement to General Intelligence}

A key limitation of current DGM-style systems is that they are commonly evaluated through task-specific benchmarks. While benchmark improvement is necessary, it is not sufficient for general intelligence.

We define general intelligence in this context as:

\begin{quote}
The ability to acquire, represent, abstract, transfer, and apply knowledge across changing tasks and domains under limited resources.
\end{quote}

This definition emphasizes five capabilities:

\begin{enumerate}[leftmargin=*]
    \item \textbf{Acquisition}: learning from experience.
    \item \textbf{Representation}: storing knowledge in reusable form.
    \item \textbf{Abstraction}: extracting general patterns from specific experiences.
    \item \textbf{Transfer}: applying knowledge across domains.
    \item \textbf{Resource allocation}: focusing limited compute and attention effectively.
\end{enumerate}

A DGM system that only stores agent code and benchmark scores lacks the structure needed to support these capabilities. It may discover useful modifications, but it does not necessarily convert experience into reusable knowledge.

EvoMind-DGM addresses this by introducing a persistent space-backed knowledge substrate and by using MOSES, ECAN, and PLN to evolve, prioritize, and reason over knowledge-bearing cognitive structures.

\section{EvoMind-DGM Architecture}

\subsection{Overview}

EvoMind-DGM is a neuro-symbolic architecture for open-ended evolution of self-improving agents. Conceptually, it still combines DGM-style evolution, persistent symbolic memory, symbolic policy evolution, attention allocation, probabilistic reasoning, and LLM-based mutation. Implementation-wise, however, these are treated as cooperating subsystems coordinated by a DGM evolution controller rather than as one already-integrated monolithic runtime. Its high-level loop is:

\begin{enumerate}[leftmargin=*]
    \item A parent agent or policy is selected from the archive.
    \item Relevant prior knowledge is retrieved from the space-backed memory layer.
    \item The symbolic runtime invokes PLN to reason about likely causes, risks, and transfer opportunities.
    \item The symbolic runtime invokes ECAN to allocate attention to promising agents, memories, tasks, and hypotheses.
    \item The symbolic runtime invokes MOSES to evolve symbolic cognitive policies or procedures.
    \item An LLM proposes implementation-level modifications or task-specific strategies.
    \item A child agent is constructed.
    \item The child agent is evaluated across multiple domains.
    \item Results, traces, and extracted abstractions are stored in the space-backed memory layer.
    \item Archive entries, attention values, and probabilistic beliefs are updated.
\end{enumerate}

The architecture can be summarized as:

\begin{verbatim}
DGM Evolution Controller
    |                    \
    |                     \
    v                      v
LLM Mutation       MeTTa/PeTTa Symbolic Runtime
                           |
                           +-- Space-Backed Memory
                           +-- PLN Reasoning
                           +-- ECAN Attention
                           +-- MOSES Policy Evolution
              \            /
               v          v
                 Child Agent
                      |
                      v
             Multi-Domain Evaluation
                      |
                      v
     Archive and Symbolic Memory Update
\end{verbatim}

\subsection{Core Components}

\subsubsection{Agent Archive}

The archive stores agent variants, symbolic policies, lineages, evaluation histories, and metadata. Each archive entry may include:

\begin{verbatim}
AgentVariant {
    id
    parent_id
    code_or_configuration
    symbolic_policy
    task_scores
    transfer_scores
    safety_scores
    attention_value
    mutation_description
    causal_hypotheses
    evaluation_logs
}
\end{verbatim}

Unlike the original DGM archive, EvoMind-DGM's archive is linked to the symbolic memory layer, allowing agent histories to become part of the system's long-term knowledge.

\subsubsection{Space-Backed Knowledge Substrate}

The knowledge substrate stores structured knowledge extracted from agent experiences. In a current implementation, this layer may be realized using Hyperon spaces, PeTTa-managed MeTTa spaces, or MORK-backed spaces while remaining conceptually Atomspace-inspired. This includes:

\begin{itemize}[leftmargin=*]
    \item episodic knowledge: what happened during a task,
    \item semantic knowledge: general concepts and task structures,
    \item procedural knowledge: reusable policies and workflows,
    \item meta-cognitive knowledge: knowledge about the system's own reasoning,
    \item goal/value knowledge: constraints and optimization objectives.
\end{itemize}

An example representation may be:

\begin{verbatim}
(AgentVariant Agent-17)
(ParentOf Agent-9 Agent-17)
(UsesPolicy Agent-17 ReflectBeforeToolUse)
(Improves Agent-17 ResearchComprehension)
(TransfersTo ReflectBeforeToolUse CodingTask)
(TruthValue (TransfersTo ReflectBeforeToolUse CodingTask) 0.72 0.61)
\end{verbatim}

\subsubsection{MOSES Symbolic Evolution Layer}

MOSES searches over symbolic policies represented in a constrained grammar. A policy may take features of the task, agent state, memory state, and uncertainty estimates as input, and output a cognitive action. In implementation terms, MOSES candidates may be generated and evaluated inside a MeTTa-centered symbolic runtime and then exposed to the controller as reusable policies.

Example policy schema:

\begin{verbatim}
IF uncertainty > theta_1 AND task_difficulty = high
THEN call_verifier

IF memory_relevance > theta_2 AND task_type = research
THEN retrieve_memory

IF progress < theta_3 AND tool_failures > 1
THEN replan
\end{verbatim}

MOSES evolves such policies based on multi-objective fitness, including task success, transfer, interpretability, safety, and compute cost.

\subsubsection{ECAN Attention Layer}

ECAN assigns attention values to entities in the symbolic memory layer and the archive. A simplified attention score may be:

\begin{equation}
Att(x) =
\alpha I(x) +
\beta N(x) +
\gamma U(x) +
\delta T(x) +
\eta S(x) -
\lambda C(x)
\end{equation}

where:

\begin{itemize}[leftmargin=*]
    \item $I(x)$ is recent improvement potential,
    \item $N(x)$ is novelty,
    \item $U(x)$ is uncertainty,
    \item $T(x)$ is transfer potential,
    \item $S(x)$ is safety relevance,
    \item $C(x)$ is evaluation cost.
\end{itemize}

ECAN guides parent selection, task selection, memory retrieval, evaluation depth, and safety prioritization. In a practical implementation, its scores are made available to the DGM controller as part of the symbolic runtime state.

\subsubsection{PLN Reasoning Layer}

PLN reasons over symbolic-memory knowledge to estimate causal and transfer relationships. In the current ecosystem, this layer can be naturally implemented in PeTTa/MeTTa and queried by the DGM controller during mutation and evaluation planning.

Example inference targets include:

\begin{verbatim}
Policy P improves task family F.
Policy P transfers from domain D1 to domain D2.
Mutation M caused improvement I.
Agent A has high risk under task type T.
Task T will reduce uncertainty about hypothesis H.
\end{verbatim}

PLN truth values can be used by ECAN and MOSES. For instance, if PLN infers that a policy has high transfer potential but low confidence, ECAN may allocate more attention to testing that policy.

\subsubsection{LLM Mutation Layer}

LLMs provide semantic flexibility. They may:

\begin{itemize}[leftmargin=*]
    \item generate code modifications,
    \item translate symbolic policies into executable agent behavior,
    \item summarize task failures,
    \item propose new tasks,
    \item explain mutation rationales,
    \item generate hypotheses for PLN to evaluate,
    \item convert natural-language experience into symbolic-memory structures.
\end{itemize}

The LLM is therefore not the whole cognitive system. It is a generative and interpretive component within a larger neuro-symbolic loop.

\section{Knowledge-Guided Self-Improvement}

The key architectural difference between DGM and EvoMind-DGM is the role of knowledge.

In standard DGM, the loop is:

\begin{equation}
\text{Agent} \rightarrow \text{Mutation} \rightarrow \text{Evaluation} \rightarrow \text{Archive}
\end{equation}

In EvoMind-DGM, the loop becomes:

\begin{equation}
\text{Experience} \rightarrow \text{Knowledge} \rightarrow \text{Reasoning} \rightarrow \text{Attention} \rightarrow \text{Evolution} \rightarrow \text{New Experience}
\end{equation}

This allows the system to accumulate knowledge about its own learning process.

\subsection{Knowledge Extraction}

After each evaluation, the system extracts knowledge from the child agent's behavior. This may include:

\begin{itemize}[leftmargin=*]
    \item task success or failure,
    \item tools used,
    \item reasoning steps,
    \item memory retrieval quality,
    \item verifier usage,
    \item failure causes,
    \item policy decisions,
    \item safety events,
    \item compute cost,
    \item cross-domain effects.
\end{itemize}

This information is converted into Atomspace representations.

\subsection{Abstraction and Transfer}

The system should not merely memorize task outcomes. It should abstract reusable patterns.

For example, from multiple task experiences, it may infer:

\begin{verbatim}
Reflect-before-tool-use improves performance when uncertainty is high.
Evidence mapping improves research-comprehension tasks.
Verifier calls improve math accuracy but increase cost.
Memory retrieval helps long-context tasks but can harm short tasks.
\end{verbatim}

These abstractions can then guide future mutations and evaluations.

\section{Research Questions}

This work investigates the following research questions:

\begin{enumerate}[leftmargin=*]
    \item \textbf{RQ1}: Can a DGM-style self-improving agent accumulate reusable knowledge in Atomspace that improves future learning and adaptation?
    \item \textbf{RQ2}: Can MOSES-evolved symbolic cognitive policies transfer across task domains better than purely LLM-generated prompt or code modifications?
    \item \textbf{RQ3}: Can ECAN-based attention allocation reduce evaluation cost while preserving or improving open-ended exploration?
    \item \textbf{RQ4}: Can PLN-based reasoning over agent histories improve mutation selection, task selection, and safety assessment?
    \item \textbf{RQ5}: Does EvoMind-DGM produce more interpretable and generalizable improvements than DGM-only baselines?
\end{enumerate}

\section{Experimental Design}

\subsection{Systems Compared}

We propose comparing four systems:

\begin{enumerate}[leftmargin=*]
    \item \textbf{Baseline LLM Agent}: fixed workflow, no self-improvement.
    \item \textbf{DGM-only Agent}: archive-based self-modification using LLM-generated mutations.
    \item \textbf{DGM + MOSES + ECAN}: symbolic policy evolution and attention-guided resource allocation.
    \item \textbf{EvoMind-DGM}: full knowledge-guided architecture with Atomspace, MOSES, ECAN, PLN, and LLM mutation.
\end{enumerate}

\subsection{Task Domains}

To evaluate generality, tasks should span multiple domains:

\begin{itemize}[leftmargin=*]
    \item software engineering and code repair,
    \item mathematical or logical reasoning,
    \item research-paper comprehension,
    \item tool-use and planning tasks,
    \item memory-dependent tasks,
    \item safety and robustness tests.
\end{itemize}

The goal is not merely to maximize performance in each domain independently, but to test whether improvements discovered in one domain transfer to others.

\subsection{Evolvable Objects}

EvoMind-DGM may evolve:

\begin{itemize}[leftmargin=*]
    \item reflection policies,
    \item verifier-selection policies,
    \item memory-retrieval policies,
    \item task-decomposition procedures,
    \item tool-use policies,
    \item archive-selection rules,
    \item fitness aggregation functions,
    \item task-selection policies,
    \item safety-check policies.
\end{itemize}

\subsection{Metrics}

Evaluation should include both task performance and general-intelligence-oriented metrics.

\begin{table}[h]
\centering
\begin{tabular}{ll}
\toprule
Metric & Description \\
\midrule
Task success & Performance on individual task domains \\
Transfer score & Improvement transferred across domains \\
Knowledge reuse & Frequency and utility of retrieved Atomspace knowledge \\
Policy reuse & Number of domains benefiting from the same symbolic policy \\
Improvement per compute & Performance gain per LLM call or evaluation cost \\
Archive diversity & Diversity of retained agent/policy variants \\
Interpretability & Human readability of evolved policies and hypotheses \\
Safety stability & Absence of safety regressions during self-improvement \\
Adaptation speed & Ability to improve on new task families with fewer examples \\
\bottomrule
\end{tabular}
\caption{Proposed evaluation metrics for EvoMind-DGM.}
\end{table}

\section{Expected Contributions}

This paper makes the following expected contributions:

\begin{enumerate}[leftmargin=*]
    \item A conceptual architecture extending Darwin G\"odel Machines from task-specific agent improvement toward knowledge-guided general intelligence.
    \item A proposal for using Atomspace as a persistent knowledge substrate for self-improving agents.
    \item A mechanism for evolving reusable symbolic cognitive policies using MOSES.
    \item An ECAN-based approach to attention and resource allocation in open-ended agent evolution.
    \item A PLN-based reasoning layer for inferring causality, transfer, risk, and expected improvement.
    \item An experimental framework for evaluating generality through knowledge reuse, policy transfer, and adaptation efficiency.
\end{enumerate}

\section{Safety Considerations}

Self-improving agents introduce significant safety concerns. EvoMind-DGM addresses these concerns by making safety part of the architecture rather than treating it as an external filter.

Safety-relevant knowledge is stored in Atomspace, including:

\begin{itemize}[leftmargin=*]
    \item unsafe mutations,
    \item tool misuse events,
    \item hallucination patterns,
    \item failed verification cases,
    \item policies associated with risk,
    \item constraints on self-modification.
\end{itemize}

MOSES can evolve explicit safety policies. ECAN can allocate attention to risky or uncertain branches. PLN can reason about the probability that a mutation creates a safety regression. Archive selection can reject or quarantine agents that improve capability while reducing safety.

A simplified archive decision rule may be:

\begin{verbatim}
IF safety_regression(agent) > threshold:
    reject_or_quarantine(agent)
ELSE IF transfer_gain(agent) > 0 AND capability_gain(agent) > 0:
    archive(agent)
ELSE IF novelty(agent) high AND risk(agent) low:
    archive_as_stepping_stone(agent)
\end{verbatim}

\section{Limitations}

Several limitations remain.

First, implementing a full EvoMind-DGM system requires orchestration across an external DGM controller, LLM agents, MeTTa/PeTTa runtimes, and space implementations rather than simply enabling a single integrated Hyperon package. Second, evaluating general intelligence remains difficult, and benchmark design must avoid narrow overfitting. Third, translating unstructured LLM traces into reliable symbolic-memory representations is non-trivial. Fourth, MOSES-evolved policies may be interpretable but still difficult to validate across open-ended environments. Fifth, ECAN and PLN must be adapted to operate effectively over large and evolving agent archives. Sixth, version skew and interface mismatch across Hyperon-family repositories may introduce substantial engineering overhead.

This paper therefore presents EvoMind-DGM as a research architecture and experimental agenda rather than a completed AGI system.

\section{Conclusion}

Darwin G\"odel Machines provide a powerful framework for open-ended self-improvement, but task-specific benchmark improvement alone is insufficient for general intelligence. General intelligence requires systems that accumulate knowledge, abstract from experience, transfer strategies across domains, reason under uncertainty, and allocate attention under limited resources.

EvoMind-DGM extends DGM with a Hyperon-inspired neuro-symbolic architecture centered on a DGM controller, a MeTTa/PeTTa symbolic runtime, and persistent space-backed memory. Within this architecture, MOSES evolves reusable symbolic cognition, ECAN allocates attention and compute, PLN reasons over uncertainty and transfer, and LLMs provide semantic and implementation-level generation. Together, these components define a path from self-improving agents toward knowledge-guided general intelligence.

The central research challenge is to determine whether such systems can acquire reusable cognitive knowledge over time and use it to improve adaptation across domains. If successful, EvoMind-DGM would represent a step toward open-ended, interpretable, and knowledge-bearing self-improving intelligence.

\bibliographystyle{plain}
\bibliography{references}

\end{document}
```
